% 第 4 章 IC Validator Dashboard
\bichapter{IC Validator Dashboard}{IC Validator Dashboard}
\label{chap:dashboard}

本章介绍 IC Validator Dashboard 的功能。

IC Validator Dashboard 包括以下各节：
\begin{itemize}
  \item 简介
  \item 从 IC Validator 获取 Dashboard 电子邮件更新
  \item 独立启动 IC Validator Dashboard
  \item 初始目录设置
  \item IC Validator 图形用户界面
\end{itemize}

% ============================================================
\section{简介}
\subsection*{Introduction}

IC Validator Dashboard 可方便地查看运行进展、规则总数与已完成数、已报告的违例数量，以及 IC Validator 如何使用已分配的资源。

芯片临近完成时，物理验证运行处于流片关键路径上，确保按时完成愈发重要。可将 IC Validator Dashboard 与 IC Validator 的 \cmd{-host\_elastic}、\cmd{-host\_add} 和 \cmd{-host\_remove} 命令行选项配合使用，以优化物理验证运行与硬件资源。

% ============================================================
\section{从 IC Validator 获取 Dashboard 电子邮件更新}
\subsection*{Getting IC Validator Dashboard Email Updates From IC Validator}

IC Validator Dashboard 可通过电子邮件发送更新。IC Validator 根据写入 \cmd{icv.log} 文件的已完成更新百分比自动发送。启动运行时在 IC Validator 命令行指定 \cmd{-email dashboard} 即可启用。例如：

\begin{lstlisting}
icv ... -email dashboard ... runset.rs
\end{lstlisting}

邮件包含三页内容，其中 IC Validator Dashboard 显示按内存使用着色的 MRI 快照，以及违例密度热图。

\figplaceholder{Figure 7: IC Validator Dashboard Email Contents}{IC Validator Dashboard 电子邮件内容}{fig:dashboard-email}

% ============================================================
\section{独立启动 IC Validator Dashboard}
\subsection*{Starting IC Validator Dashboard Standalone}

可用下列命令启动 IC Validator Dashboard：

\begin{lstlisting}
icv_dashboard [-email ["user1,user2,..."]] [-V] [-h]
 [<run_details_path>|<dp_log_file>]
\end{lstlisting}

若不提供任何参数，Dashboard 会在 \cmd{./run\_details} 目录中查找分布式处理日志文件，并打开图形用户界面（GUI）。

命令行选项见表~\ref{tab:dashboard-cli}。

\begin{longtable}{@{}>{\raggedright\arraybackslash\ttfamily}p{0.38\textwidth} p{0.54\textwidth}@{}}
\caption{IC Validator Dashboard 命令行选项}\label{tab:dashboard-cli}\\
\toprule
\textbf{\textrm{选项}} & \textbf{\textrm{说明}} \\
\midrule
\endfirsthead
\caption[]{IC Validator Dashboard 命令行选项（续）}\\
\toprule
\textbf{\textrm{选项}} & \textbf{\textrm{说明}} \\
\midrule
\endhead
\bottomrule
\endfoot
-email ["user1,user2,\ldots"] & 向用户发送一次电子邮件更新后退出。 \\
-V & 打印产品版本后退出。忽略其他命令行选项。 \\
-h & 打印程序选项后退出。忽略其他命令行选项。 \\
run\_details\_path$|$dp\_log\_file & 使用指定的 run details 目录或 DP 日志文件显示 GUI。默认值为 \cmd{./run\_details/}。指定 run details 目录时，使用该目录中的分布式处理日志文件。若该目录中存在多个分布式处理日志文件，则必须显式指定要使用的日志文件。 \\
\end{longtable}

\begin{noteBox}
以此方式启动 IC Validator Dashboard 时，请确保 \cmd{PATH} 变量包含 \cmd{\$ICV\_HOME\_DIR/bin}，并且 \cmd{LM\_LICENSE\_FILE} 或 \cmd{SNPSLMD\_LICENSE\_FILE} 环境变量指向带有 IC Validator license 的有效 Synopsys license 服务器。
\end{noteBox}

% ============================================================
\section{初始目录设置}
\subsection*{Initial Directory Setting}

IC Validator Dashboard 会自动创建初始目录以存放中间文件。可通过设置环境变量 \cmd{DASHBOARD\_TMP} 指定初始目录路径。若该路径已设置且有效，则 Dashboard 在此路径下创建初始目录；否则依次尝试当前目录、用户主目录。若无法创建初始目录，将显示错误。

% ============================================================
\section{IC Validator 图形用户界面}
\subsection*{The IC Validator Graphical User Interface}

使用上一节所述命令行选项交互式启动 GUI 后，它会持续监视所找到或所指向的 IC Validator 运行。GUI 打开时显示 Summary 页面；从运行加载信息后，显示内容如图~\ref{fig:dashboard-summary}。

\subsection{IC Validator Dashboard Summary 页面}
\subsubsection*{IC Validator Dashboard Summary Page}

Summary 页面包含五个面板，分别提供状态、CPU 使用、内存使用、运行进度与规则进度信息。将指针悬停在规则进度或状态“指示灯”上直至出现 tooltip，可查看更多信息。

\figplaceholder{Figure 8: IC Validator Dashboard Summary Page}{IC Validator Dashboard Summary 页面}{fig:dashboard-summary}

\figplaceholder{Figure 9: IC Validator Dashboard Status Panel}{IC Validator Dashboard 状态面板}{fig:dashboard-status-panel}

\subsubsection{IC Validator Dashboard 状态}
\paragraph*{IC Validator Dashboard Status}

IC Validator Dashboard 有五种状态值：
\begin{itemize}
  \item \textbf{Waiting}：Dashboard GUI 已启动，正在等待 IC Validator 运行开始以及分布式处理日志文件出现。
  \item \textbf{Running}：IC Validator 正在运行，Dashboard 正在等待更新。
  \item \textbf{Completed}：IC Validator 运行已成功完成。
  \item \textbf{Stopped}：IC Validator 运行因错误结束。
  \item \textbf{Unknown}：Dashboard 已读入不完整的分布式处理日志文件，但无法判断作业是否仍在运行。例如：加载了部分 DP 日志文件，而 IC Validator 已不再运行。
\end{itemize}

\subsubsection{IC Validator 运行状态指示器}
\paragraph*{IC Validator Run Status Indicators}

通常这些状态指示器为绿色。当 IC Validator 磁盘空间或内存不足时，相应指示器变为红色。若命令崩溃、主机被丢弃，或总运行时间中有相当大部分被延迟命令（deferred commands）占用，则 ``Other'' 状态指示器变为红色。将指针悬停在指示器上可显示状态原因。

\subsubsection{IC Validator 运行信息}
\paragraph*{IC Validator Run Information}

状态面板的这一部分显示 IC Validator 的主要输入。

\subsubsection{CPU 使用面板}
\paragraph*{CPU Usage Panel}

CPU 使用面板显示该 IC Validator 运行在所有主机上使用的 CPU 总数。以 elastic 模式运行时，作业运行期间可能释放部分 CPU；此时峰值使用量有可能超过已分配的 CPU 数。

\figplaceholder{Figure 10: IC Validator Dashboard CPU Usage Panel}{IC Validator Dashboard CPU 使用面板}{fig:dashboard-cpu-panel}

\subsubsection{内存使用面板}
\paragraph*{Memory Usage Panel}

Memory Usage 面板显示当前内存使用最多的主机的内存使用情况。当为 IC Validator 运行分配了多台主机时，该主机在运行过程中可能发生变化。

\figplaceholder{Figure 11: IC Validator Dashboard Memory Usage Panel}{IC Validator Dashboard 内存使用面板}{fig:dashboard-mem-panel}

\subsubsection{运行进度面板}
\paragraph*{Run Progress Panel}

Run Progress 面板显示自运行开始以来的时长（elapsed time）以及 IC Validator 已完成操作的百分比。进度条以图形方式表示完成百分比。

\figplaceholder{Figure 12: IC Validator Dashboard Run Progress Panel}{IC Validator Dashboard 运行进度面板}{fig:dashboard-run-progress}

\subsubsection{规则进度面板}
\paragraph*{Rule Progress Panel}

Rule Progress 面板显示已执行的规则数以及已生成的错误数。如图~\ref{fig:dashboard-summary} 所示，将指针悬停在进度条上时，tooltip 会显示精确数值。

\figplaceholder{Figure 13: IC Validator Dashboard Rule Progress Panel}{IC Validator Dashboard 规则进度面板}{fig:dashboard-rule-progress}

在 Explorer 模式下运行 IC Validator 时，Summary 页面的 Rule Progress 面板会显示最近的 Explorer Tier 消息。将指针悬停在消息上可显示 Tooltip，以查看更详细信息。

\subsection{IC Validator Dashboard CPU History 页面}
\subsubsection*{IC Validator Dashboard CPU History Page}

如图~\ref{fig:dashboard-cpu-history} 所示，CPU 页面显示按主机着色编码的总 CPU 使用历史。上方面板是随时间变化的图：显示已分配的 CPU 总数以及各主机正在使用的 CPU 数。下方面板作为图例，并显示各主机的当前状态快照。将指针悬停在不同元素上会出现带有附加信息的 tooltip。

\figplaceholder{Figure 14: IC Validator Dashboard CPU History Page}{IC Validator Dashboard CPU History 页面}{fig:dashboard-cpu-history}

\figplaceholder{Figure 15: CPU History Page When Job Completes}{作业完成时的 CPU History 页面}{fig:dashboard-cpu-history-done}

\subsection{IC Validator Memory History 页面}
\subsubsection*{IC Validator Memory History Page}

如图~\ref{fig:dashboard-mem-history} 所示，Memory 页面显示按主机的内存使用历史。上方面板显示各主机随时间的内存使用。下方面板显示各主机内存使用的当前快照：实心条表示当前内存使用，彩色轮廓表示峰值内存使用，黑色轮廓表示主机的总物理内存。若使用 IC Validator 的 \cmd{-host\_memory} 命令行选项指定了主机内存限制，则黑色轮廓中高于该限制的部分以灰色填充。将指针悬停在不同元素上会出现带有附加信息的 tooltip。

\figplaceholder{Figure 16: The IC Validator Dashboard Memory History Page}{IC Validator Dashboard Memory History 页面}{fig:dashboard-mem-history}

\subsection{IC Validator Dashboard 选项}
\subsubsection*{IC Validator Dashboard Options}

单击菜单栏上的邮件图标可打开 \textbf{IC Validator Dashboard -- Options} 对话框。可在此选择最多三种不同的电子邮件通知选项。基于时间的邮件在 GUI 打开期间每 10 分钟发送一次。基于事件的邮件在每次打印新进度行时发送。若 IC Validator 以 \cmd{-email dashboard} 命令行选项启动，则该框为勾选状态。取消勾选可禁用基于事件的邮件；也可在进行中的 ICV 作业完成或失败时发送通知邮件。

\figplaceholder{Figure 17: IC Validator Dashboard - Options Dialog Box}{IC Validator Dashboard -- Options 对话框}{fig:dashboard-options}

Dashboard 电子邮件包含 GUI 各选项卡所显示信息的快照。
